% od-core.tex : shared base of the optodraw 1D symbol library.
% Colors, line styles, the /od/ key family, the port macro and the read-back
% macros. Family files (od-onchip.tex, ...) hold one pic per component and are
% loaded after this file. The library needs TikZ and nothing else.
\usetikzlibrary{calc,arrows.meta}

% Line art is black. Color carries a fixed meaning and nothing else.
\definecolor{odheat}{HTML}{F0E442}                         % heater metal
\definecolor{odp}{HTML}{E8615F}                            % p doping (light red)
\definecolor{odn}{HTML}{56B4E9}                            % n doping (light blue)
\definecolor{odw1}{HTML}{0072B2}\definecolor{odw2}{HTML}{009E73} % wavelengths
\definecolor{odw3}{HTML}{E69F00}\definecolor{odw4}{HTML}{CC79A7}

\newif\ifodshowports \odshowportsfalse
% The look. 1d (the default): line art, color only where it means something.
% 2d: the same drawing painted, with materials and depth (gradients, a cast
% shadow). The 2D art turns it on with \odtwodlook in the scope of a part.
\def\odlook{1d}\def\odlooktwod{2d}
\newcommand\odtwodlook{\def\odlook{2d}\definecolor{odtrace}{HTML}{FB923C}\color{ods700}}  % a trace glows; lines are slate
% od paint={<1D options>}{<2D options>}: one style, two looks.
\tikzset{od paint/.code 2 args={\ifx\odlook\odlooktwod\pgfkeysalso{#2}\else\pgfkeysalso{#1}\fi}}
% Materials in 2D: the slate greys of modern interface design, light and muted.
% Each is the top and the bottom of a soft gradient, and a hairline edge in a
% darker tint of the same colour. No black line.
\definecolor{ods50}{HTML}{F8FAFC}\definecolor{ods100}{HTML}{F1F5F9}\definecolor{ods200}{HTML}{E2E8F0}
\definecolor{ods300}{HTML}{CBD5E1}\definecolor{ods400}{HTML}{94A3B8}\definecolor{ods500}{HTML}{64748B}
\definecolor{ods600}{HTML}{475569}\definecolor{ods700}{HTML}{334155}\definecolor{ods800}{HTML}{1E293B}
\definecolor{ods900}{HTML}{0F172A}
\colorlet{odcaseA}{white}\colorlet{odcaseB}{ods200}\colorlet{odcaseE}{ods400}       % instrument case
\colorlet{odanoA}{ods300}\colorlet{odanoB}{ods400}\colorlet{odanoE}{ods500}         % a housing (light, not black)
\colorlet{odaluA}{ods50}\colorlet{odaluB}{ods300}\colorlet{odaluE}{ods400}          % aluminium, steel
\definecolor{odglassA}{HTML}{F4FBFF}\definecolor{odglassB}{HTML}{CDE7F4}\definecolor{odglassE}{HTML}{7FB2CF} % glass
\definecolor{odgoldA}{HTML}{FBE7A1}\definecolor{odgoldB}{HTML}{D4AF37}\definecolor{odgoldE}{HTML}{A8841F}  % gold, brass
\definecolor{odpcbA}{HTML}{34D399}\definecolor{odpcbB}{HTML}{059669}\definecolor{odpcbE}{HTML}{047857}    % circuit board
\colorlet{odscreen}{ods800}\colorlet{odedge}{ods500}                % a screen; the edge of a painted part
\definecolor{odcopperA}{HTML}{EDC2A0}\definecolor{odcopperB}{HTML}{B87333}\definecolor{odcopperE}{HTML}{7A4A1E} % copper, as the 3D parts
% The light falls from above. A shading maps the box of a path onto its middle, 25bp to 75bp.
\pgfdeclareverticalshading{odbox}{100bp}{color(0bp)=(ods500); color(25bp)=(ods500); color(31bp)=(ods400); color(55bp)=(ods300); color(70bp)=(ods200); color(75bp)=(ods100); color(100bp)=(ods100)}   % a housing: a soft top edge
\pgfdeclareverticalshading{odcylinder}{100bp}{color(0bp)=(ods500); color(25bp)=(ods500); color(36bp)=(ods400); color(50bp)=(ods200); color(62bp)=(ods50); color(69bp)=(ods200); color(75bp)=(ods400); color(100bp)=(ods400)}   % a tube, its highlight above the middle
\pgfdeclarehorizontalshading{odrod}{100bp}{color(0bp)=(ods400); color(25bp)=(ods400); color(38bp)=(ods200); color(46bp)=(ods50); color(58bp)=(ods200); color(75bp)=(ods500); color(100bp)=(ods500)}   % a post: light from the left
\pgfdeclareverticalshading{odsilver}{100bp}{color(0bp)=(ods300); color(25bp)=(ods300); color(45bp)=(ods100); color(58bp)=(white); color(66bp)=(ods100); color(75bp)=(ods300); color(100bp)=(ods300)}   % a mirror face
\pgfdeclareverticalshading{odblack}{100bp}{color(0bp)=(ods900); color(25bp)=(ods900); color(60bp)=(ods800); color(71bp)=(ods700); color(75bp)=(ods600); color(100bp)=(ods600)}   % matte black anodising: a sheen on the top edge only
\definecolor{odbeam}{HTML}{EF4444}                                  % a beam drawn inside a part, in 2D
% Line scale: every line width is multiplied by it. 1 in the symbols; the 2D art
% sets it per part (optodraw/symbols/art.py) so lines shrink with a small part.
\def\odls{1}

% Keys that more than one family uses. A family file adds its own keys.
\pgfkeys{/od/.cd,
  length/.initial=2,      % waveguide, ring bus, interferometer arm length
  color/.initial=black,   % resonator color (a wavelength, in a weight bank)
  inputs/.initial=1,
  outputs/.initial=2,
  pitch/.initial=0.5,     % port pitch
  kind/.initial=,         % which variant of a pic: see the pic (phase, converter)
  ports/.initial=,        % which set of port names: see the pic
  count/.initial=8,       % arrays and banks: how many
  view/.initial=,         % top: the view from above, for a part whose 1D symbol is a side view (2D art)
  sample/.initial=,       % no: a holder drawn without its own sample (a cryostat a sample is placed in)
  sync/.initial=,         % yes: a laser with a sync output (a jack under it, an electrical port)
  show ports/.is if=odshowports,
}
\newcommand\odget[1]{\pgfkeysvalueof{/od/#1}}

\tikzset{
  od wg/.style={od paint={line width=0.9pt*\odls}{line width=0.55pt*\odls, draw=ods600}, line cap=round, line join=round},   % optical
  od ring/.style={line width=1.7pt*\odls},                                  % the resonator itself
  od el/.style={od paint={line width=0.55pt*\odls}{line width=0.4pt*\odls, draw=ods600}, line cap=round, line join=round},  % electrical
  od ray/.style={od paint={line width=0.35pt*\odls}{line width=0.3pt*\odls, draw=ods500}, line cap=round},
  od band/.style={draw=black!80, line width=0.4pt*\odls, line join=round},  % a strip of metal or doping
  od ground/.style={od paint={od band, fill=white}{od hair, draw=odgoldE, top color=odgoldA, bottom color=odgoldA!60!odgoldB}},
  od signal/.style={od paint={od band, fill=black!22}{od hair, draw=odgoldE, top color=odgoldA, bottom color=odgoldB}},
  % 2D materials. In 1D each is the plain line art it was before.
  % 2D: a hairline edge, the line every painted part uses.
  od hair/.style={line width=0.3pt*\odls, line cap=round, line join=round},
  % A soft cast shadow: two translucent copies of the outline, low and wide (no library needed).
  od shadow/.style={preaction={draw=none, fill=ods900, fill opacity=0.03, transform canvas={xshift=0.012cm, yshift=-0.06cm}},
                    preaction={draw=none, fill=ods900, fill opacity=0.04, transform canvas={xshift=0.008cm, yshift=-0.042cm}},
                    preaction={draw=none, fill=ods900, fill opacity=0.05, transform canvas={xshift=0.004cm, yshift=-0.026cm}},
                    preaction={draw=none, fill=ods900, fill opacity=0.06, transform canvas={yshift=-0.012cm}}},
  od shine/.style={},                                                % the shadings hold the light now
  od case/.style={od paint={od wg, fill=white}{od hair, draw=odcaseE, top color=odcaseA, bottom color=odcaseB, od shadow}},
  od housing/.style={od paint={od wg, fill=white}{od hair, draw=odanoE, shading=odbox, od shadow}},
  od tube/.style={od paint={od wg, fill=white}{od hair, draw=odanoE, shading=odcylinder, od shadow}},
  od alu/.style={od paint={od wg, fill=white}{od hair, draw=odaluE, top color=odaluA, bottom color=odaluB, od shadow}},
  od alusmall/.style={od paint={od el, fill=white}{od hair, draw=odaluE, top color=odaluA, bottom color=odaluB}},
  od pcb/.style={od paint={od wg, fill=white}{od hair, draw=odpcbE, top color=odpcbA, bottom color=odpcbB, od shadow}},
  % What stops light, and the dark head of a laser: black in 1D, matte black anodising in 2D.
  od absorber/.style={od paint={od el, fill=black}{od hair, draw=ods900, shading=odblack}},
  od rod/.style={od paint={od el, fill=white}{od hair, draw=odanoE, shading=odrod, shading angle=0}},     % a post: steel
  od copper/.style={od paint={od el, fill=black!22}{od hair, draw=odcopperE, top color=odcopperA, bottom color=odcopperB}},
  od chipcase/.style={od paint={od wg, fill=white}{od hair, draw=ods800, top color=ods600, bottom color=ods800}},
  od pixel/.style={od paint={od ray, fill=white}{od hair, draw=ods500, top color=ods600, bottom color=ods700}},
  od knob/.style={od paint={od el, fill=white}{od hair, draw=ods400, top color=white, bottom color=ods300, od shadow}},
  od jack/.style={od paint={od el, fill=white}{od hair, draw=ods500, top color=ods100, bottom color=ods400}},
  od port/.style={circle, draw, fill=white, line width=0.5pt*\odls, inner sep=0pt, minimum size=3.4pt},
  od eport/.style={rectangle, draw, fill=white, line width=0.5pt*\odls, inner sep=0pt, minimum size=3.0pt},
  od light/.style={od el, -{Stealth[length=3.2pt,width=2.6pt]}},
  od cable/.style={od el, draw=black!65},                             % a cable between parts: thin, grey, solid
}
\def\odmw{0.12}   % half width of every metal strip: electrodes, heaters, segments

% A port: \odport[<kind>]{<name>}{<position>}. The kind is optodraw's port kind
% (optical, electrical, thermal, fiber, free_space); the default is optical.
% The kind hub is not a port: it marks where the beams of a splitter meet, and a
% view draws the beam from each port that a scene uses to it.
% This gives the coordinate (<pic name>-<name>), an optional mark (circle =
% light, square = electrical or thermal), and an entry in the list that the
% build reads back.
\gdef\odportlist{}
\newcommand\odport[3][optical]{%
  \coordinate (-#2) at (#3);
  \ifodshowports
    \def\odtmpa{#1}\def\odtmpb{electrical}\def\odtmpc{thermal}\def\odtmpd{hub}%
    \ifx\odtmpa\odtmpb \node[od eport] at (#3) {};%
    \else\ifx\odtmpa\odtmpc \node[od eport] at (#3) {};%
    \else\ifx\odtmpa\odtmpd
    \else \node[od port] at (#3) {};\fi\fi\fi
  \fi
  \xdef\odportlist{\odportlist\ifx\odportlist\empty\else,\fi#2/#1}%
}

% Read-back. The build compiles one pic named (s) at the origin and reads its
% ports and bounding box from <job>.ports.
\newwrite\odout
\newcommand\odwriteports[1]{%
  \foreach \odpname/\odpkind in \odportlist {%
    \path (#1-\odpname); \pgfgetlastxy{\odpx}{\odpy}%
    \immediate\write\odout{port \odpname\space\odpkind\space\odpx\space\odpy}}%
}
\newcommand\odwritebbox{%
  \path (current bounding box.south west); \pgfgetlastxy{\odpx}{\odpy}%
  \immediate\write\odout{sw \odpx\space\odpy}%
  \path (current bounding box.north east); \pgfgetlastxy{\odpx}{\odpy}%
  \immediate\write\odout{ne \odpx\space\odpy}%
}
